% !TEX root = % !TEX root = /Users/nai1ka/Projects/PerfX/Thesis/thesis.tex
\chapter{Literature Review}
\label{chap:lr}
\chaptermark{Literature Review}

This chapter reviews existing research on mobile application profiling and monitoring. It looks at how performance metrics are collected and at which metrics matter most for mobile development. It then examines the main approaches to performance regression detection, which prepares the choice of method made later in this work.

The review is organised into three main parts. The first part compares the two main approaches to collecting performance metrics in mobile applications. The second part summarises the metrics used in previous research and explains how they contribute to the assessment of application performance and user experience. The third part reviews existing approaches to performance regression detection, from manual threshold-based alerting to release-to-release comparison in large-scale production systems. Together, these parts show the current state of research and provide a foundation for identifying the gap that the proposed system aims to address.

\section{Performance monitoring}

Performance strongly shapes how users experience a mobile application and whether they keep using it. Slow response times, high energy use, or a laggy interface frustrate users and often end in poor reviews and uninstalls~\cite{Khalid2015}. Catching such problems early keeps an application usable and its users satisfied, as Tang \emph{et al.} \cite{Tang} demonstrate. Two strategies are commonly used to understand application performance: profiling and monitoring. Profiling runs during development or testing and yields detailed, high-resolution data. Monitoring, in contrast, takes place after the application is released and runs on real
user devices. It must introduce as little overhead as possible to prevent any interference with normal app usage.

The Android ecosystem adds complexity to performance evaluation due to fragmentation - the wide variety of hardware,
operating system versions, and user behaviours \cite{Fragmentation}. This diversity makes it difficult to reproduce rare
or device-specific issues found in the wild. The same app may behave differently on various devices depending on their
environment, OS version, and hardware specifications \cite{Zikan}. This problem motivates research into tools and
methods for collecting and analysing mobile performance data effectively.

\section{Approaches for Performance Data Collection}
\subsection{Pre-release Profiling}

Several tools have been developed to support performance profiling during development. One example is
ProfileDroid by Wei \emph{et al.} \cite{ProfileDroid}, which uses a multi-layer approach to profiling. It gathers data from four layers: application metadata, user-interface interaction traces, operating-system metrics such as CPU and memory usage, and network activity. Looking at these layers together lets developers see how the parts of an application interact, for instance how a single UI event triggers a CPU spike or a network request. The main strength of ProfileDroid is
its ability to correlate performance across layers, giving a complete system-level view. However, its high overhead and
dependence on a controlled environment limit its practical use on real devices.

Another relevant project is AndroStep \cite{AndroStep}, which focuses on analysing Android storage performance during
I/O operations. The authors found that existing Android benchmarking tools did not simulate realistic app behaviour,
especially missing operations like fsync() followed by 4 KB random writes - a common and essential operation used by
SQLite, the embedded relational database built into Android for local data storage \cite{AndroStep}. The tool is useful for benchmarking Android storage, but it measures nothing beyond storage and has to be run by hand, which limits it to a test environment rather than a live one.

% TODO i think i need to add soemthing else here
PowDroid, introduced by Bouaffar \emph{et al.} in 2021 \cite{PowDroid}, was designed to measure the energy consumption
of applications. However, like most pre-release tools, it works only under controlled conditions, which makes it
difficult to test applications on a large range of real devices.

In the field of UI performance, a new state-of-the-art work, GuiWatcher, was proposed by Liu \emph{et al.}
\cite{GuiWatcher}. They presented an advanced way to detect UI lags and freezes. GUIWatcher uses computer vision
techniques on screen recordings and is able to detect three types of problematic frames: "junky" frames (dropped or
delayed ones), long loading frames (lengthy transition), and frozen frames. This method reflects what users actually see
rather than relying only on timing logs. However, it is computationally expensive and suitable only for use during
development or continuous integration testing. It also raises privacy concerns, because it records the device screen.

ADEL, by Zhang \emph{et al.} \cite{ADEL}, targets network performance and energy leaks. It extends the platform to tag downloaded data and follow how that data moves through an application, which reveals unnecessary network activity that wastes energy. The tool is not built for continuous monitoring, however, because it modifies the Android system and adds high overhead, so it cannot run on real user devices.


% TODO add somewhere: As noted by Xia et al. (2022), enabling full method tracing in Android Studio can slow apps down
% by over 10×, making it “unsuitable for continuous use in production”



\subsection{Post-release Monitoring}

The main idea of post-release profiling is to collect metrics from apps as they run on real user devices. The main
advantage of this approach is that it captures realistic usage patterns and real-world diversity that are hard to
reproduce in a controlled test environment. The main challenge is to minimise overhead to avoid affecting user experience.

One of the first studies in this area was AppInsight \cite{AppInsight}, which introduced a method for instrumenting
applications without modifying their source code. The framework automatically injects instrumentation into the compiled
bytecode and collects runtime traces. By collecting traces from 30 users over several months, the authors were able to
identify performance bottlenecks that were not visible during development. However, this approach is limited because
developers cannot easily add custom metrics.

Another example is Hubble \cite{Hubble}, which is currently integrated into all Huawei devices. The main feature of
Hubble is that it monitors the performance in an efficient in-memory buffer, where older data is constantly overwritten
\cite{Hubble}. The relevant metrics are only persisted when a performance issue is detected, providing engineers with
detailed runtime traces from just before it occurred. This approach combines high efficiency with low storage cost.

PerfProbe \cite{PerfProbe} works across layers to find the root cause of a performance problem rather than only its symptom. It traces both the application and the system on real devices at low overhead, then uses statistical analysis to flag functions that behave abnormally and tie them to resource factors such as CPU, memory, or I/O. A developer thus learns not only where a slowdown happens but why. Across 22 Android apps, PerfProbe diagnosed real issues and reduced user-perceived latency in six of them by 32\% to 86\% \cite{PerfProbe}.

Panappticon \cite{Panappticon} is a lightweight tracing system that records performance across the application, system, and kernel layers. It traces critical execution paths automatically by linking each user input to the display update it produces. Unlike AppInsight, it can also catch problems that arise from interactions between several apps or system components. Over a one-month study its overhead stayed around 6\%, low enough for continuous use in the field.


\section{Performance Metrics}
Performance is not captured by a single number but by a set of metrics, which fall into two kinds: those that describe how an application uses device resources, and those that describe what the user actually perceives. Resource metrics explain where a performance problem originates, while user-perceived metrics show how strongly it affects the user. The two are linked, because resource pressure usually surfaces as a slower or less responsive interface.

On the resource side, the most common metrics are CPU usage, memory usage, storage I/O, and energy consumption. High CPU utilisation overloads threads and reduces responsiveness, and CPU contention has been shown to be a frequent cause of poor user experience on Android \cite{Yaohui}. Memory usage matters for stability: an application that consumes too much RAM can trigger the system's low-memory killer or frequent garbage-collection pauses, so developers profile heap usage to catch allocation bottlenecks and memory leaks \cite{DroidPerf}. % TODO random citation - check
Storage I/O also affects latency, mostly on devices with slower flash memory, where a slow read can make a screen take noticeably longer to open \cite{AndroStep}. Energy use is just as visible. Surveys put battery life at the top of users' performance concerns, and people avoid applications they consider too power-hungry \cite{Saraiva}, so finding energy leaks is an important part of mobile optimisation.

On the user-perceived side, several more specific metrics describe responsiveness. Frame rendering time is the most direct of these. A frame has to be drawn within the display's refresh budget, roughly 16.7 ms on a 60 Hz screen, and a frame that misses this budget shows up as stutter or a dropped frame \cite{android_slow_rendering}. Frames per second (FPS) carries the same information in a different form and follows directly from frame time. Startup time, the wait before an application becomes usable after launch, largely sets the user's first impression \cite{AndroidStartup}. % TODO add citation: Android app startup time (vitals/launch-time)
Interaction latency is the delay between a user action and the first visible response. Even when CPU and memory look normal, a long delay here noticeably reduces satisfaction, and guidance such as Google’s RAIL model treats a response slower than about one second as the point where the user’s attention drifts \cite{PerfProbe}. The most severe responsiveness failure is an Application Not Responding (ANR) event, raised when the main thread stays blocked for several seconds, which often makes the user abandon the application \cite{AndroidANR}.

In practice no single metric is sufficient, and tools combine several to build a fuller picture. ProfileDroid and Panappticon correlate UI and system activity, while PerfProbe links CPU and I/O data with execution paths. The choice of metrics depends on the goal: pre-release profiling collects many low-level metrics, whereas post-release monitoring favours a smaller set of user-perceived metrics, such as frame time, startup time, and interaction latency, that are cheap to measure on real devices and reflect more directly what the user actually experiences.

\section{Performance Regression Detection}

A performance regression is a situation where software still works correctly but performs worse than a previous baseline, for example slower response times or higher resource usage. In mobile applications this usually appears after an update, as increased frame rendering time, slower startup, higher CPU usage, or greater memory consumption. Regression detection aims to identify a sustained degradation in performance across a population of users between two releases, rather than isolated, short-lived spikes.

\subsection{Approaches to Detecting Performance Regressions}

The simplest and most widely used approach in industry is threshold-based alerting. Tools such as Firebase Performance Monitoring \cite{FirebasePerf} and Google Play Vitals \cite{GooglePlayVitals} let developers set fixed limits on individual metrics, for example a maximum startup time, and raise an alert when a limit is exceeded. This approach is easy to understand, but the thresholds must be set manually for each metric and application, and a single fixed value rarely fits every device and usage pattern, which leads to false alerts or missed regressions. Such tools also react to absolute values rather than to changes between versions, so a slowdown introduced by a new release goes unnoticed as long as it stays within the configured limit.

A more adaptive approach compares current performance against a historical baseline instead of a fixed limit. Nguyen \emph{et al.} \cite{Nguyen2012} build control limits from a baseline run and flag a new run when it deviates too far from them, which removes the need to choose absolute thresholds by hand. Their method, however, was designed for controlled load-test runs of server software rather than for data collected continuously from real user devices. On mobile, Carat \cite{Carat} applies a similar idea at the population level: it collects energy data from many users, builds a baseline of normal behaviour for each application, and reports devices and apps that deviate from it. Carat compares devices rather than consecutive releases, so it can locate an abnormal app but does not attribute a regression to a particular version. Industry observability platforms such as Datadog go further and compare live metrics against an automatically learned historical or seasonal baseline, flagging statistically significant deviations \cite{Datadog}. These tools react to anomalies as they unfold over time rather than to the difference between two releases, and they target server and infrastructure metrics rather than mobile applications on user devices.

The approach closest to the needs of this work compares each new release directly against the previous one, so that a detected regression is attributed to a specific version. This is the principle behind canary analysis, which is widely used in industry to gate deployments by comparing a new version against a stable baseline before full rollout \cite{Kayenta}.
At large scale, Meta's FBDetect \cite{FBDetect} monitors hundreds of thousands of time series in production, compares versions to catch even small regressions, and filters out false alerts caused by transient issues such as load spikes or rolling updates. These systems show that release-to-release comparison is an effective and interpretable way to detect regressions, because each alert is tied to a concrete change. Both systems run inside the data centre, where the software executes on controlled and largely homogeneous servers. They therefore do not deal with device fragmentation and have not been applied to performance data collected from heterogeneous Android devices after release.


% TODO maybe remove
\subsection{Mobile-Specific Considerations and Research Gap}

Applying these ideas to Android raises additional constraints. Devices differ widely in hardware and operating system version, so the same release can perform differently across devices, and comparing them directly can make normal hardware differences look like regressions. Devices are also resource-constrained, and collecting detailed traces from real users raises privacy concerns. A practical design therefore keeps the on-device part lightweight and moves the heavier analysis to a server.

Despite progress in both server-side regression detection and mobile monitoring, no existing solution combines continuous post-release monitoring on Android devices with automated, server-side detection of regressions between releases. Industry monitoring tools either provide no automated alerting or rely on fixed absolute thresholds that must be set manually for each metric, and the release-to-release comparison used in large-scale systems has not been adapted to Android, where device heterogeneity has to be taken into account. The proposed system addresses this gap by comparing consecutive releases on the server and grouping the comparison by device performance cohort.


\section{Summary}
This chapter reviewed existing research on mobile performance profiling, monitoring, and regression detection. Pre-release tools such as ProfileDroid, AndroStep and PowDroid offer deep system-level insights but introduce high overhead and require controlled environments, making them unsuitable for real user devices. Post-release monitoring systems such as AppInsight, PerfProbe, and Hubble avoid these drawbacks by gathering data straight from real users at minimal overhead.

The chapter also separated two kinds of metric: resource use such as CPU, memory, storage I/O, and energy, and the user-perceived responsiveness that those resources ultimately affect. It closed with regression detection, moving from manual threshold-based alerting to baseline and release-to-release comparison in large-scale systems such as FBDetect. The recurring finding is that comparing a new release against the one before it is an effective way to catch sustained degradation.

Even so, no existing solution brings together continuous, low-overhead post-release monitoring and automated server-side regression detection built for Android. This gap is what the proposed system sets out to fill.